\kchapter{Kubernetes Architecture}{control-plane}

\begin{objectives}
\begin{itemize}
  \item what a cluster, a node, and the control plane are;
  \item the job of every control-plane component: API server, etcd,
        scheduler, controller manager, cloud controller manager;
  \item the job of every node component: kubelet, kube-proxy, container runtime;
  \item what actually happens, step by step, when you create a Deployment.
\end{itemize}
\end{objectives}

\section{The cluster at a glance}

A Kubernetes \kterm{cluster} is a set of machines (physical or virtual) called
\kterm{nodes}, split into two roles:

\begin{itemize}
  \item the \kterm{control plane} --- the ``brain'' that stores desired state and
        makes decisions;
  \item \kterm{worker nodes} --- the ``muscle'' that actually runs your containers.
\end{itemize}

\begin{center}
\resizebox{\linewidth}{!}{%
\begin{tikzpicture}[
  comp/.style={draw=k8sblue, fill=white, rounded corners=4pt, line width=0.8pt,
               minimum width=2.75cm, minimum height=2.2cm, align=center,
               font=\footnotesize\bfseries, text=k8snavy},
  small/.style={draw=k8sblue, fill=white, rounded corners=4pt, line width=0.8pt,
               minimum width=2.3cm, minimum height=1.9cm, align=center,
               font=\scriptsize\bfseries, text=k8snavy},
  pod/.style={fill=k8sblue!15, draw=k8sblue!60, rounded corners=3pt,
              minimum width=1.25cm, minimum height=0.9cm, font=\scriptsize, align=center},
  >={Stealth[length=2.6mm]}]

  % ---------- control plane
  \node[comp] (etcd)  at (-6.2,0) {\kimg[1cm]{etcd}\\etcd};
  \node[comp] (api)   at (-3.1,0) {\kimg[1cm]{api}\\kube-apiserver};
  \node[comp] (sched) at (0,0)    {\kimg[1cm]{sched}\\kube-scheduler};
  \node[comp] (cm)    at (3.1,0)  {\kimg[1cm]{c-m}\\kube-controller-\\manager};
  \node[comp] (ccm)   at (6.2,0)  {\kimg[1cm]{c-c-m}\\cloud-controller-\\manager};
  \begin{scope}[on background layer]
    \node[fill=k8slight, draw=k8sblue, line width=1.2pt, rounded corners=10pt,
          inner sep=0.45cm, fit=(etcd)(ccm), label={[font=\bfseries, text=k8sblue]above:\faBrain\ \ CONTROL PLANE}] (cp) {};
  \end{scope}
  \draw[<->, thick, k8sblue] (etcd) -- (api);
  \draw[<->, thick, k8sblue] (api) -- (sched);
  \draw[<->, thick, k8sblue] (api.north) to[bend left=22] (cm.north);
  \draw[<->, thick, k8sblue] (api.north) to[bend left=28] (ccm.north);

  % cloud
  \node[font=\Huge, text=k8sgray] (cloud) at (10.0,0) {\faCloud};
  \node[font=\scriptsize, text=k8sgray, below=-1pt of cloud] {Cloud API};
  \draw[<->, thick, k8sgray, dashed] (ccm) -- (cloud);

  % user
  \node[font=\Huge, text=k8snavy] (user) at (-10.0,0) {\faUserCircle};
  \node[font=\scriptsize\bfseries, text=k8snavy, below=0pt of user] {you + kubectl};
  \draw[->, line width=1.4pt, k8sorange] (user) -- node[above, font=\scriptsize, text=k8sorange]{HTTPS} (cp.west |- user);

  % ---------- worker nodes
  \foreach \nx/\nn in {-4.6/1, 4.6/2}{
    \node[small] (kubelet\nn) at (\nx-2.6,-5.0) {\kimg[0.9cm]{kubelet}\\kubelet};
    \node[small] (proxy\nn)   at (\nx,-5.0)     {\kimg[0.9cm]{k-proxy}\\kube-proxy};
    \node[small] (cri\nn)     at (\nx+2.6,-5.0) {\kimg[0.9cm]{containerd}\\container\\runtime};
    \node[pod] (pa\nn) at (\nx-2.6,-7.0) {\kimg[0.45cm]{pod}\\Pod};
    \node[pod] (pb\nn) at (\nx-0.9,-7.0) {\kimg[0.45cm]{pod}\\Pod};
    \node[pod] (pc\nn) at (\nx+0.8,-7.0) {\kimg[0.45cm]{pod}\\Pod};
    \node[pod] (pd\nn) at (\nx+2.6,-7.0) {\kimg[0.45cm]{pod}\\Pod};
    \begin{scope}[on background layer]
      \node[fill=k8sblue!6, draw=k8sdeep, line width=1.1pt, rounded corners=10pt,
            inner sep=0.35cm, fit=(kubelet\nn)(cri\nn)(pa\nn)(pd\nn),
            label={[font=\bfseries\small, text=k8sdeep]below:\faServer\ \ WORKER NODE \nn}] (node\nn) {};
    \end{scope}
    \draw[<->, thick, k8sblue] (kubelet\nn.north) -- (kubelet\nn.north |- cp.south);
    \draw[->, thick, k8sgray] (kubelet\nn) -- (cri\nn.west |- kubelet\nn) ;
  }
\end{tikzpicture}}
\end{center}

\begin{keyidea}
\textbf{Everything talks to the API server, and only the API server talks to
etcd.} The scheduler, controllers, kubelets, \kc{kubectl}, and your CI system
are all just clients of one HTTP API.
\end{keyidea}

\section{Control-plane components}

\subsection{\kicon{api}\ kube-apiserver --- the front door}

The API server exposes the Kubernetes REST API. Every create, read, update, or
delete --- from you, from a controller, or from a node --- goes through it. For
each request it:

\begin{enumerate}
  \item \textbf{authenticates} the caller (who are you? certificate, token, OIDC);
  \item \textbf{authorises} the action (are you allowed? usually RBAC);
  \item runs \textbf{admission control} (should this be modified or rejected?
        e.g.\ inject defaults, enforce policies);
  \item \textbf{validates} the object against its schema;
  \item \textbf{persists} it to etcd and notifies anyone \emph{watching}.
\end{enumerate}

The API server is stateless, so you can run several copies behind a load balancer.

\subsection{\kicon{etcd}\ etcd --- the memory}

etcd is a distributed, strongly consistent key--value store. It holds the
\emph{entire} state of the cluster: every Pod, Service, Secret, and config. It
uses the Raft consensus algorithm, so production clusters run 3 or 5 etcd
members and tolerate the loss of 1 or 2.

\begin{warningbox}
If you lose etcd and have no backup, you lose your cluster's definition.
Back it up regularly (Chapter~\ref{ch:production}). Managed services do this for you.
\end{warningbox}

\subsection{\kicon{sched}\ kube-scheduler --- the matchmaker}

When a new Pod has no node assigned, the scheduler picks one. It works in two phases:

\begin{itemize}
  \item \textbf{Filtering}: remove nodes that cannot run the Pod (not enough
        CPU/memory, wrong labels, taints, missing volumes, port conflicts).
  \item \textbf{Scoring}: rank the remaining nodes (spread replicas apart,
        prefer nodes with the image cached, balance utilisation) and pick the best.
\end{itemize}

The scheduler only \emph{decides}. It writes the chosen node name into the Pod
object; it does not start anything.

\subsection{\kicon{c-m}\ kube-controller-manager --- the autopilots}

A single binary running dozens of controllers, each owning one kind of object:
the Deployment controller, ReplicaSet controller, Node controller (notices dead
nodes), Job controller, EndpointSlice controller, ServiceAccount controller,
and many more. Each runs the reconcile loop from Chapter 2.

\subsection{\kicon{c-c-m}\ cloud-controller-manager --- the cloud bridge}

Only present on cloud clusters. It talks to the cloud provider's API to create
load balancers for Services, attach routes, and learn when cloud VMs are deleted.

\section{Node components}

\subsection{\kicon{kubelet}\ kubelet --- the node agent}

Runs on every node. It watches the API server for Pods assigned to its node,
asks the container runtime to start their containers, mounts volumes, runs
health probes, and reports status back. If a container dies, the kubelet
restarts it.

\subsection{\kicon{k-proxy}\ kube-proxy --- the traffic cop}

Implements Services on each node by programming packet-forwarding rules
(iptables, IPVS, or nftables) so that traffic to a Service's virtual IP lands on
a healthy Pod. Some network plugins (like Cilium) replace kube-proxy with eBPF.

\subsection{\kicon{containerd}\ Container runtime}

The software that actually pulls images and runs containers ---
\textbf{containerd} or \textbf{CRI-O}. The kubelet talks to it through the
\kterm{Container Runtime Interface} (CRI). Images built with Docker run fine,
because they are standard OCI images.

\section{Add-ons you will meet}

\begin{center}
\ktablestart
\begin{tabularx}{\linewidth}{P{3.6cm} L}
\kheader \kth{Add-on} & \kth{Purpose} \\
CoreDNS & Cluster DNS: lets Pods reach \kc{my-service.my-namespace} by name. \\
CNI plugin & Pod networking. Examples: Calico, Cilium, Flannel, kindnet. \\
CSI drivers & Storage plugins for cloud disks, NFS, Ceph, etc. \\
metrics-server & Live CPU/memory figures for \kc{kubectl top} and autoscaling. \\
Ingress / Gateway controller & HTTP(S) routing from outside the cluster. \\
\end{tabularx}
\end{center}

\section{Walkthrough: what happens when you create a Deployment}

You run \kc{kubectl apply -f web.yaml}, asking for a Deployment with 3 replicas.

\begin{center}
\resizebox{\linewidth}{!}{%
\begin{tikzpicture}[
  actor/.style={fill=k8sblue, text=white, rounded corners=4pt, minimum width=2.3cm,
                minimum height=0.9cm, font=\footnotesize\bfseries, align=center},
  msg/.style={->, >={Stealth[length=2.4mm]}, thick, k8snavy},
  lab/.style={font=\scriptsize, fill=white, inner sep=1.5pt, text=k8snavy}]
  \foreach \n/\x/\txt in {k/0/kubectl, a/3/API server, e/6/etcd, d/9/Deployment\\controller,
                        r/12/ReplicaSet\\controller, s/15/Scheduler, l/18/kubelet}{
    \node[actor] (\n) at (\x,0) {\txt};
    \draw[k8smid, dashed, thick] (\n.south) -- (\x,-9.3);
  }
  \draw[msg] (0,-1.0) -- node[lab, above]{1. POST Deployment} (3,-1.0);
  \draw[msg] (3,-1.7) -- node[lab, above]{2. store} (6,-1.7);
  \draw[msg] (3,-2.5) -- node[lab, above]{3. watch event: new Deployment} (9,-2.5);
  \draw[msg] (9,-3.3) -- node[lab, above]{4. create ReplicaSet} (3,-3.3);
  \draw[msg] (3,-4.1) -- node[lab, above]{5. watch: new ReplicaSet} (12,-4.1);
  \draw[msg] (12,-4.9) -- node[lab, above]{6. create 3 Pods (no node yet)} (3,-4.9);
  \draw[msg] (3,-5.7) -- node[lab, above]{7. watch: unscheduled Pods} (15,-5.7);
  \draw[msg] (15,-6.5) -- node[lab, above]{8. bind each Pod to a node} (3,-6.5);
  \draw[msg] (3,-7.3) -- node[lab, above]{9. watch: Pod assigned to me} (18,-7.3);
  \draw[msg] (18,-8.1) -- node[lab, above]{10. start containers, report status Running} (3,-8.1);
  \draw[msg] (3,-8.8) -- node[lab, above]{11. persist status} (6,-8.8);
\end{tikzpicture}}
\end{center}

Notice that no component orders another around. Each one simply \emph{watches}
the API for objects it cares about and reacts. This loose coupling is what makes
Kubernetes so robust and extensible.

\begin{deepdive}
Controllers do not poll. They open long-lived \emph{watch} connections to the
API server and keep a local cache (an ``informer''). Every object carries a
\kc{resourceVersion}; updates use optimistic concurrency, so two controllers can
never silently overwrite each other's changes.
\end{deepdive}

\section{Kubernetes objects: spec and status}

Almost every object you create has two halves:

\begin{itemize}
  \item \kc{spec} --- written by \emph{you}: the desired state;
  \item \kc{status} --- written by \emph{Kubernetes}: the observed state.
\end{itemize}

Controllers spend their lives trying to make \kc{status} match \kc{spec}.

\begin{recap}
\begin{itemize}
  \item A cluster = control plane (API server, etcd, scheduler, controller
        managers) + worker nodes (kubelet, kube-proxy, runtime).
  \item The API server is the only door; etcd is the only memory.
  \item Components cooperate by watching the API, not by calling each other.
  \item Objects have a \kc{spec} (desired) and a \kc{status} (actual).
\end{itemize}
\end{recap}
